\documentclass[../principal.tex]{subfiles}
\graphicspath{{\subfix{../imagenes/}}}

\begin{document}

\chapter{Arreglos}

\section{Definición}

Los arreglos son colecciones o conjuntos de datos.

Los arreglos son también conocidos en inglés como arrays.

En algunos lenguajes como Java y C, los arreglos nos permiten almacenar datos del mismo tipo, por ejemplo, solamente números enteros, o solamente caracteres, pero en JavaScript podemos almacenar distintos tipos de datos en el mismo arreglo.

\section{Sintaxis en JavaScript}

Los arreglos en JavaScript se pueden declarar con la palabra clave let, igual que las variables, y asignarle como valor inicial un arreglo vacío con corchetes.

\begin{lstlisting}[language=JavaScript]
let unArreglo = [];
\end{lstlisting}


\end{document}


